\documentclass[conference]{IEEEtran}

\usepackage[utf8]{inputenc}
\usepackage{booktabs}
\usepackage{url}
\usepackage{hyperref}

\renewcommand{\ttdefault}{cmtt}

\hypersetup{
  colorlinks=true,
  linkcolor=black,
  citecolor=black,
  urlcolor=black
}

\title{WFH Burnout Predictor: A CRISP-DM Machine Learning Project}

\author{
\IEEEauthorblockN{Simone Mastria}
\IEEEauthorblockA{Machine Learning and Data Mining, Module 2\\
University of Bologna\\
Email: simone.mastria@studenti.unibo.it}
}

\begin{document}

\maketitle

\begin{abstract}
This paper describes a machine learning project for estimating burnout in
work-from-home settings from daily behavioral indicators such as working hours,
screen time, meetings, breaks, after-hours work, sleep and task completion. The
work follows the CRISP-DM methodology discussed in the course The main task is regression
on \texttt{burnout\_score}; classification of \texttt{burnout\_risk} is treated
as a secondary analysis because the high-risk class is strongly imbalanced.
\end{abstract}

\begin{IEEEkeywords}
machine learning, CRISP-DM, burnout prediction, work from home, regression,
classification, fairness
\end{IEEEkeywords}

\section{Introduction}

The goal of the project is to develop a complete machine learning workflow for the
prediction of employee burnout in work-from-home (WFH) scenarios. Burnout is a
sensitive topic: predictions should not be interpreted as medical diagnoses or
used for automatic human-resources decisions. The intended use is decision
support at an aggregate level, for example to identify workload patterns that may
deserve organizational attention.

The project uses the public Kaggle dataset \emph{Work From Home Employee Burnout
Dataset}~\cite{kaggle}. It contains 1,800 daily observations, 180 users and 10
records per user. Each row includes workload and lifestyle variables:
\texttt{day\_type}, \texttt{work\_hours}, \texttt{screen\_time\_hours},
\texttt{meetings\_count}, \texttt{breaks\_taken},
\texttt{after\_hours\_work}, \texttt{sleep\_hours} and
\texttt{task\_completion\_rate}. The available targets are
\texttt{burnout\_score}, a continuous score, and \texttt{burnout\_risk}, a
categorical label with values \texttt{Low}, \texttt{Medium} and \texttt{High}.

The work was carried out individually by Simone Mastria and covers dataset
inspection, preprocessing design, feature engineering, model selection,
evaluation and interpretation.

The main challenge is not the implementation of a single algorithm, but the
construction of a defensible data mining workflow. In
particular, the dataset contains repeated observations for each employee. If a
standard random row split were used, the same employee could appear in both
training and test data, producing an overly optimistic evaluation. Another
challenge is the semantic interpretation of the available variables: some
features are plausible behavioral predictors, while others may be close proxies
for the target itself. These aspects motivated a conservative methodology based
on grouped splitting, baseline comparison and explicit discussion of limitations.

\section{Proposed Method}

\subsection{CRISP-DM Strategy}

The project was organized according to CRISP-DM~\cite{shearer}. In the business
understanding phase, the objective was formulated as a predictive and
interpretable support tool for WFH burnout monitoring. In the data understanding
phase, the dataset schema, missing values, duplicates, target distributions and
correlations were analyzed. In the data preparation phase, transformations were
implemented through scikit-learn pipelines to avoid leakage. Modeling compared
baselines and supervised learning algorithms. Evaluation considered both
technical metrics and domain limitations. Deployment was described as a cautious
plan rather than an automatic production system.

Table~\ref{tab:crispdm} summarizes how the CRISP-DM phases were mapped to the
project artifacts. This mapping was used both to structure the notebook and to
prepare the final paper.

\begin{table}[h]
\centering
\caption{CRISP-DM mapping for the project.}
\label{tab:crispdm}
\begin{tabular}{p{0.26\linewidth}p{0.62\linewidth}}
\toprule
Phase & Project artifact \\
\midrule
Business understanding & Goal, target choice, metrics and risks \\
Data understanding & Schema, target distribution, correlations and bias notes \\
Data preparation & Grouped split, encoding, scaling and feature engineering \\
Modeling & Baselines, candidate models and tuning \\
Evaluation & Test metrics, error analysis and interpretability \\
Deployment & Colab execution, saved pipeline and monitoring plan \\
\bottomrule
\end{tabular}
\end{table}

\subsection{Data Preparation}

The dataset has no missing values and no duplicated rows. However, it still
requires preprocessing, The identifier
\texttt{user\_id} was not used as a feature because it is not generalizable.
Instead, it was used as a grouping variable during splitting: records from the
same user must not appear in both training and test sets.

Table~\ref{tab:variables} reports the core variables. The distinction between
predictive features, grouping variable and targets is essential to avoid leakage.

\begin{table}[h]
\centering
\caption{Main variables used in the project.}
\label{tab:variables}
\begin{tabular}{ll}
\toprule
Variable & Role \\
\midrule
\texttt{user\_id} & Grouping variable, excluded from predictors \\
\texttt{day\_type} & Categorical predictor \\
\texttt{work\_hours} & Numeric predictor \\
\texttt{screen\_time\_hours} & Numeric predictor \\
\texttt{meetings\_count} & Numeric predictor \\
\texttt{breaks\_taken} & Numeric predictor \\
\texttt{after\_hours\_work} & Binary predictor \\
\texttt{sleep\_hours} & Numeric predictor \\
\texttt{task\_completion\_rate} & Numeric predictor, high-leverage feature \\
\texttt{burnout\_score} & Regression target \\
\texttt{burnout\_risk} & Classification target \\
\bottomrule
\end{tabular}
\end{table}

The train/test split was performed with \texttt{GroupShuffleSplit}, using 80\%
of the records for training and 20\% for testing. This produced 1,440 training
rows and 360 test rows, with zero overlap between train and test users. This is
more realistic than a random row split because the model is evaluated on unseen
users.

Categorical variables were encoded with one-hot encoding and numeric variables
were standardized with \texttt{StandardScaler}. All transformations were placed
inside a \texttt{ColumnTransformer} and a scikit-learn \texttt{Pipeline}; thus,
encoders and scalers are fitted only on training data.

Several interpretable engineered features were added:
\texttt{work\_intensity}, \texttt{screen\_work\_ratio},
\texttt{meeting\_density}, \texttt{sleep\_deficit} and
\texttt{overwork\_flag}. These variables express workload pressure, screen
exposure, meeting concentration and sleep deficit.

The preprocessing pipeline also improves reproducibility. Every candidate model
receives exactly the same transformed feature matrix, and the same logic is
available at inference time if the fitted pipeline is saved. This is preferable
to manually transforming training and test sets in separate notebook cells,
which would make leakage and inconsistent preprocessing more likely.

\subsection{Modeling Strategy}

The primary task is regression on \texttt{burnout\_score}. The secondary task is
classification on \texttt{burnout\_risk}. This choice is important because the
categorical target is highly imbalanced: 1,527 rows are \texttt{Low}, 253 are
\texttt{Medium}, and only 20 are \texttt{High}. Therefore, a classifier can
obtain high accuracy while failing on the most critical cases.

For regression, the following models were evaluated: \texttt{DummyRegressor},
Ridge Regression, K-Nearest Neighbors, Random Forest and Hist Gradient Boosting.
For classification, the models were \texttt{DummyClassifier}, balanced Logistic
Regression, balanced Random Forest and Hist Gradient Boosting. Cross-validation
used \texttt{GroupKFold} on the training set. A light hyperparameter search was
then performed for Random Forest regression with \texttt{RandomizedSearchCV}.

The main regression metric is mean absolute error (MAE), supported by RMSE and
$R^2$. The classification metrics are macro F1, balanced accuracy and
class-level precision/recall. These metrics are aligned with the course
recommendation to define clear, task-dependent evaluation criteria rather than
optimizing accuracy blindly.

The selected algorithms cover different inductive biases. Ridge Regression is a
linear baseline with regularization and is useful to check whether the target is
mostly explained by linear relationships. KNN is sensitive to scaling and local
neighborhood structure. Random Forest and Hist Gradient Boosting can capture
non-linear effects and interactions. Dummy models provide lower bounds: if a
candidate does not clearly beat them, it should not be considered useful.

Hyperparameter optimization was deliberately limited. The course material on
AutoML emphasizes that automated search is useful, but it should not replace
problem understanding. Therefore, only a small randomized search was applied to
Random Forest after comparing the initial candidates. The final test set was not
used to choose hyperparameters.

\section{Results}

\subsection{Data Insights}

The data understanding phase showed that \texttt{task\_completion\_rate} is the
dominant feature: it has a correlation of about $-0.96$ with
\texttt{burnout\_score}. This is predictive, but it must be interpreted with
caution. It may be a direct proxy for burnout or part of the synthetic logic used
to generate the target. Other variables such as \texttt{work\_hours},
\texttt{screen\_time\_hours} and \texttt{meetings\_count} have weaker positive
correlations with burnout.

The descriptive analysis also confirmed that the dataset is clean in a
technical sense: there are no missing values, no duplicated rows, and each user
has exactly 10 observations. However, technical cleanliness is not the same as
data quality for decision making. Since the dataset is public and its collection
context is limited, representativeness cannot be guaranteed. This limits the
strength of any claim about real organizations.

An explicit data quality report was added to the notebook, including checks for
invalid ranges, duplicates and missing values. Statistical outliers were also
detected through the IQR method for the main numeric variables. They were not
removed automatically: in this domain, extreme burnout scores or extreme working
patterns may be the most relevant observations rather than errors.

The risk-label distribution is a major issue. The \texttt{High} class represents
only about 1.1\% of the dataset. This affects model selection, evaluation and
fairness: the most important class from a preventive perspective is also the
least represented.

For this reason, the project treats \texttt{burnout\_risk} as a secondary task.
The continuous score gives a more stable learning signal, while the classes are
useful to communicate results and inspect whether high-risk observations are
missed.

\subsection{Regression Results}

Table~\ref{tab:regression} reports the test-set results for the main regression
models. All meaningful models outperform the dummy baseline by a large margin.
Ridge obtains the lowest raw test MAE among the initial models, while Random
Forest has the best cross-validation MAE and was selected for tuning.

\begin{table}[h]
\centering
\caption{Regression performance on the test set.}
\label{tab:regression}
\begin{tabular}{lccc}
\toprule
Model & MAE & RMSE & $R^2$ \\
\midrule
Ridge & 4.79 & 6.31 & 0.934 \\
Random Forest & 4.85 & 6.41 & 0.932 \\
Hist Gradient Boosting & 5.14 & 6.73 & 0.925 \\
KNN & 9.26 & 11.65 & 0.776 \\
Dummy median & 20.35 & 25.80 & -0.099 \\
\bottomrule
\end{tabular}
\end{table}

After tuning, Random Forest used 200 estimators, maximum depth 4,
\texttt{min\_samples\_leaf=8} and \texttt{max\_features=1.0}. The tuned model
obtained CV MAE 4.85, test MAE 4.88, test RMSE 6.21 and test $R^2=0.936$. This
is a strong improvement over the dummy baseline and confirms that the selected
features contain useful predictive information.

Cross-validation on the training set supported the same conclusion. Random
Forest obtained mean CV MAE 4.97, Ridge 5.06, Hist Gradient Boosting 5.19, KNN
9.03 and the dummy baseline 18.46. The gap between training and test error was
monitored to detect overfitting: tree-based models obtained lower training error
but did not collapse on unseen users, while KNN remained substantially weaker.

The error analysis in the notebook includes a predicted-versus-actual plot, the
distribution of residuals and the largest individual errors. These checks are
important because an average error can hide systematic failures. In this domain,
large underestimation for high-burnout records would be more problematic than
small errors around low-risk observations.

To address a key methodological risk, a sensitivity analysis was performed by
removing \texttt{task\_completion\_rate}. Results dropped sharply: Ridge reached
MAE 20.25 and Random Forest MAE 20.26, close to the dummy baseline. This confirms
that \texttt{task\_completion\_rate} carries most of the predictive signal. The
model is therefore useful as a predictive pipeline, but it should not be
presented as evidence of causal burnout mechanisms.

\subsection{Classification Results}

The best classifier was balanced Logistic Regression, with balanced accuracy
0.803 and macro F1 0.646. Table~\ref{tab:classification} reports its class-level
metrics. The model performs very well on \texttt{Low} and reasonably on
\texttt{Medium}; the \texttt{High} class remains difficult because only four
high-risk examples are present in the test set.

\begin{table}[h]
\centering
\caption{Classification report for balanced Logistic Regression.}
\label{tab:classification}
\begin{tabular}{lcccc}
\toprule
Class & Precision & Recall & F1 & Support \\
\midrule
Low & 0.99 & 0.93 & 0.96 & 294 \\
Medium & 0.68 & 0.73 & 0.70 & 62 \\
High & 0.17 & 0.75 & 0.27 & 4 \\
\bottomrule
\end{tabular}
\end{table}

The low precision on \texttt{High} indicates many false alarms, while the recall
of 0.75 suggests that most observed high-risk test cases were detected. In a
preventive setting, this trade-off may be acceptable only if predictions are
reviewed by humans and not used automatically.

The classification results also show why accuracy alone would be misleading.
The dummy classifier that always predicts the majority class is not useful, even
though most observations are \texttt{Low}. Macro F1 and balanced accuracy give a
more honest view because they weight minority classes more appropriately.

\subsection{Interpretability}

Permutation importance was computed on the regression test set. The most
important feature is \texttt{task\_completion\_rate}. This is consistent with
the correlation analysis, but it should not be interpreted causally. A low task
completion rate may be a consequence, a symptom, a proxy, or a construction
component of burnout score. Other features, such as workload and screen time,
provide weaker but more directly actionable information. This distinction is
important for business interpretation: interventions should not simply pressure
employees to complete more tasks, because that may worsen the underlying
problem.

\subsection{Fairness and Ethics}

The dataset does not include sensitive attributes such as gender, age or ethnic
origin. Nevertheless, fairness cannot be ignored. The notebook checks regression
errors and high-risk recall across operational subgroups such as
\texttt{day\_type}, \texttt{after\_hours\_work} and work-hour bands. These
checks are not a full fairness audit, but they are useful to detect systematic
performance differences.

From an ethical standpoint, the safest deployment scenario is aggregate
monitoring. For example, the model could help identify teams or periods where
workload patterns become riskier. Individual predictions should require human
review, transparency and consent. The model should never be used as a punitive
tool or as an automatic performance evaluation system.

\section{Conclusions}

The project implements the required end-to-end machine learning workflow:
problem definition, public dataset selection, preprocessing, modeling,
evaluation, insight extraction and deployment discussion. The regression task is
the most reliable formulation: the tuned Random Forest reaches MAE 4.88 and
$R^2=0.936$ on unseen users, clearly outperforming the baseline. The
classification task is useful for operational interpretation, but the
\texttt{High} class is too rare for strong claims.

The most important insight is that \texttt{task\_completion\_rate} dominates the
prediction. This improves performance but also creates a methodological risk:
the variable may be too close to the target-generation process. Future work
should validate the approach on an external real-world dataset, collect richer
contextual information, define intervention thresholds with domain experts and
run a stronger fairness assessment before any deployment.


\begin{thebibliography}{1}

\bibitem{kaggle}
S. Shinde, ``Work From Home Employee Burnout Dataset,'' Kaggle. [Online].
Available:
\url{https://www.kaggle.com/datasets/sonalshinde123/work-from-home-employee-burnout-dataset}

\bibitem{shearer}
C. Shearer, ``The CRISP-DM model: The new blueprint for data mining,''
\emph{Journal of Data Warehousing}, vol. 5, no. 4, pp. 13--22, 2000.

\bibitem{francia}
M. Francia, ``Machine Learning and Data Mining, Module 2,'' University of
Bologna, course slides, A.Y. 2025/2026.

\end{thebibliography}

\end{document}
